% Minimal LaTeX template for a oneside report about kinds (Haskell, JavaScript, F#)
% No report content is added yet — section placeholders only.
\documentclass[12pt,oneside,a4paper]{article}

% Encoding and fonts
\usepackage[utf8]{inputenc}
\usepackage[T1]{fontenc}
\usepackage{lmodern}
\usepackage{microtype}

% Page geometry
\usepackage[left=25mm,right=25mm,top=25mm,bottom=25mm]{geometry}

% Mathematics and symbols
\usepackage{amsmath,amssymb}

% Graphics and floats
\usepackage{graphicx}
\usepackage{caption}

% Listings for code (light setup)
\usepackage{listings}
\usepackage{xcolor}
\lstset{basicstyle=\ttfamily\small,breaklines=true,columns=fullflexible}

% Tables
\usepackage{booktabs}

% Better lists
\usepackage{enumitem}

% Bibliography (uses biber)
\usepackage[backend=biber,style=authoryear]{biblatex}
\addbibresource{references.bib}

% Hyperlinks
\usepackage[hidelinks]{hyperref}

% Begin document
\begin{document}

% Manual title block (avoid literal "\\title" token)
\begin{center}
{\LARGE\bfseries Kinds in Haskell and Their Analogues in F\# and JavaScript}\\[1ex]
{\small Jaanis Fehling \quad --- \quad Universita di Pisa}\\[2ex]
\end{center}

\begin{abstract}
% TODO: Add a short abstract here.
\end{abstract}

% Use \csname to safely invoke the table-of-contents command
\csname tableofcontents\endcsname
\clearpage



\section{Kinds in Haskell}
The functional programming language Haskell offers besides a classical typesystem a "typesystem of a typesystem", also called \textit{kinds} (i.e. kind of a type or type of a type). In GHCI, the interactive haskell interpreter, we can use the ":t" command to check the type of a previously declared variable:
\begin{lstlisting}
    ghci> x = 42 :: Int
    ghci> :t x
    x :: Int
\end{lstlisting}
Now, the type "Int" has the following kind, which we can check using the ":k" command:
\begin{lstlisting}
    ghci> :k Int
    Int :: *
\end{lstlisting}
In contrast, we can check the kind of a type constructor such as "Maybe":
\begin{lstlisting}
    ghci> :k Maybe
    Maybe :: * -> *
\end{lstlisting}


\section{Kinds and types in JavaScript}
% TODO
\section{Kinds and types in F\#}
% TODO

\section{Conclusion}



\appendix
\section{Supplementary Material}

\printbibliography

\end{document}

